% -----------------------------------------------------------
\section{Bosom}
% -----------------------------------------------------------
	\label{BOSOM}
	
% % ----------------------
% \subsection{See also}
% % ----------------------

% ----------------------
\subsection{1828 American Dictionary of the English Language} % *1828
% ----------------------
	\label{BOSOM-1828}

	\begin{compactitem}
		\item[1] The breast of a human being and the parts adjacent.
		\item[2] The folds or covering of clothes about the breast.
		\item[3] Embrace, as with the arms; inclosure; compass; often implying friendship or affection; as, to live in the \textit{bosom} of a church.
		\item[4] The breast, as inclosing the heart; or the interior of the breast, considered as the seat of the passions.
		\item[5] The breast, or its interior, considered as a close place, the receptacle of secrets.
		\item[6] Any inclosed place; the interior; as the \textit{bosom} of the earth or of the deep.
		\item[7] The tender affections; kindness; favor; as the son of his \textit{bosom}; the wife of thy \textit{bosom}.
		\item[8] The arms, or embrace of the arms.
		\item[9] Inclination; desire. [Not used.]
		\item \textit{bosom}, in composition, implies intimacy, affection and confidence; as a \textit{bosom}-friend, an intimate or confidential friend; \textit{bosom}-lover, \textit{bosom}-interest, \textit{bosom}-secret, etc. In such phrases, \textit{bosom} may be considered as an attribute equivalent to intimate, confidential, dear.
	\end{compactitem}

% % ----------------------
% \subsection{Oxford English Dictionary} % *OED
% % ----------------------
	% \label{BOSOM-OED}

	% \begin{compactitem}
		% \item[]
	% \end{compactitem}
	
% % ----------------------
% \subsection{Restoration Lexicon} % *RL *RESTORATION LEXICON
% % ----------------------
	% \label{BOSOM-OED}

	% \begin{compactitem}
		% \item[]
	% \end{compactitem}

% ----------------------
\subsection{Scriptural Usage} % *SCRIPTURES
% ----------------------
	\label{BOSOM-SCRIPTURES}

	% % -----
	% \subsubsection{Old Testament} % *OT
	% % -----
		% \label{BOSOM-OT}
	
		% % --
		% \paragraph{KJV}
		% % --
			% \label{BOSOM-OT-KJV}
	
			% \begin{tabular}{ll}
				% Total occurrences in the OT & 
			% \end{tabular}
			
			% \begin{compactdesc}
				% \item[]
			% \end{compactdesc}
			
		% % --
		% \paragraph{Hebrew Bible}
		% % --
			% \label{BOSOM-HEBREW}
		
			% %
			% \subparagraph{\texorpdfstring{\he{sample word}}{}} % paste actual hebrew text here
			% %
				% \label{PAGA} % whatever is in step bible without periods
			
				% \begin{compactdesc}
					% \item[HALOT , PDF ] \textbf{} % Use the 1994 PDF
						% \begin{compactitem}
							% \item[1]
						% \end{compactitem}
					% \item[BDB , PDF ] \textbf{}
						% \begin{compactitem}
							% \item[1]
						% \end{compactitem}
					% \item[DCH PDF ] \textbf{} % don't include the actual page number, they repeat from volume to volume
						% \begin{compactitem}
							% \item[1]
						% \end{compactitem}
				% \end{compactdesc}
				
				% \begin{compactdesc}
					% \item[Some OT uses] \textbf{}
						% \begin{compactdesc}
							% \item[]
						% \end{compactdesc}
				% \end{compactdesc}
			
		% % --
		% \paragraph{Septuagint}
		% % --
			% \label{BOSOM-LXX}
		
			% %
			% \subparagraph{\texorpdfstring{\gr{sample word}}{}}
			% %
		
	% -----
	\subsubsection{New Testament} % *NT
	% -----
		\label{BOSOM-NT}
	
		% % --
		% \paragraph{KJV}
		% % --
			% \label{BOSOM-NT-KJV}		
	
			% \begin{tabular}{ll}
				% Total occurrences in the NT & 
			% \end{tabular}
			
			% \begin{compactdesc}
				% \item[]
			% \end{compactdesc}
			
		% --
		\paragraph{Greek New Testament} % *GREEK
		% --
			\label{BOSOM-GREEK}
		
			%
			\subparagraph{\texorpdfstring{\gr{κόλπος}}{}}
			%
				\label{KOLPOS}
			
				\begin{compactdesc}
					\item[BDAG 493a] various meanings in general literary usage, frequently with suggestion of curvature and the hollow so formed, as of a person's chest, folds in a garment or a bay of the sea; our literature contains no application of the term to anatomical parts uniquely female.
						\begin{compactitem}
							\item[1] \textbf{\textit{bosom, breast, chest}}...furthermore, apart from the idea of dining together on the same couch, ``being in someone's bosom'' denotes the closest association
								\begin{compactitem}
									\item This entry glosses \hyperref[JOHN-1-18]{John 1:18} as ``\textit{who rests in the bosom of the Father}.''
								\end{compactitem} 
						\end{compactitem}
					\item[CGL 819a, PDF 859] \textbf{}
						\begin{compactitem}
							\item[1] (singular and plural) fold at the lap of a garment (formed where it loops over the belt), \textbf{lap-fold}
							\item[2] (singular and plural) area of the body covered by the lap-fold, \textbf{lap, bosom} (usually of a mother or nurse, where a child or animal is nurtured or sheltered)
							\item[3] \textbf{bosom} (of a beautiful woman)
							\item[4] \textbf{womb} (of a woman)
							\item[5] \textbf{bosom} (of the earth, especially as holding the dead)
							\item[6] capacious depths, \textbf{bosom} (of the sea)
							\item[7] curving recess of the coast, \textbf{bay, gulf}
							\item[8] gulf of water between enclosing coasts, \textbf{gulf}
							\item[9] \textbf{valley, glen}
							\item[10] capacious hollow, \textbf{hold} (of a ship); \textbf{recess, cavity} (in the ground, in a cave)
							\item[11] \textbf{curving hollow} (of a hunting net)
						\end{compactitem}
					\item[LSJ 974a, PDF 1045] \textbf{}
						\begin{compactitem}
							\item[I.1] \textit{bosom, lap}
							\item[I.2.a] \textit{vagina}
							\item[I.2.b] supposed \textit{sinuses} in the womb
							\item[I.2.c] in poets more vaguely of the whole \textit{sinus genitalis, womb}
							\item[I.2.d] of other cavities
							\item[II] \textit{fold of a garment}
							\item[III] any \textit{bosom-like hollow}:
							\item[III.1] of the sea, first in a half-literal sense, of a sea-goddess
							\item[III.2] \textit{bay, gulf}
							\item[III.3] \textit{vale}
							\item[III.4] of a fortified site, \textit{salient}
							\item[III.5] \textit{bottom} of the chariot
							\item[III.6] \textit{fistulous ulcer} which spreads under the skin
							\item[IV] in Tactics, \textit{enveloping force}
						\end{compactitem}
				\end{compactdesc}
				
				% \begin{tabular}{ll}
					% Total occurrences in the NT & 
				% \end{tabular}
				
				% \begin{compactdesc}
					% \item[Some NT uses] \textbf{}
						% \begin{compactdesc}
							% \item[]
						% \end{compactdesc}
				% \end{compactdesc}
				
			%
			\subparagraph{TDNT: \texorpdfstring{\gr{κόλπος}}{}  (3:824--826)}
			%
				\label{KOLPOS-TDNT}
				
				``At a meal the place of the guest of honour, hence figuratively to express an inward relationship'' (824).
				
				``In the NT we find \gr{κόλπος} as a `bosom,' `lap' in \hyperref[JOHN-13-23]{John 13:23}: `at supper he took the place of honour on Jesus' breast'. Applied figuratively to membership of the community, we find the same motif in [2 Clement 4:5] in an unknown saying of the Lord...[the Greek follows; see AFGTET 142 for the Greek, and AFGTET 143 for this translation: ``For this reason, if you do these things, the Lord has said, `If you are gathered with me close to my breast (\gr{εν τω κολπω μου}), yet you do not keep my commandments, I will throw you out and say to you: ``Get away from me; I do not know where you are from, you evildoers (\gr{εργαται ανομιας})''.''']. Without the idea of a meal it expresses closest fellowship in \hyperref[JOHN-1-18]{John 1:18}.''
		
	% % -----
	% \subsubsection{Book of Mormon} % *BOM
	% % -----
		% \label{BOSOM-BOM}
	
		% \begin{tabular}{ll}
			% Total occurrences in the BOM & 
		% \end{tabular}
		
		% \begin{compactdesc}
			% \item[]
		% \end{compactdesc}
		
	% -----
	\subsubsection{Doctrine and Covenants} % *DC
	% -----
		\label{BOSOM-DC}
	
		% \begin{tabular}{ll}
			% Total occurrences in the DC & 
		% \end{tabular}
		
		\begin{compactdesc}
			\item[{\hyperref[DC109-4]{DC 109:4}}]
		\end{compactdesc}
		
	% % -----
	% \subsubsection{Pearl of Great Price} % *PGP
	% % -----
		% \label{BOSOM-PGP}
	
		% \begin{tabular}{ll}
			% Total occurrences in the PGP & 
		% \end{tabular}
		
		% \begin{compactdesc}
			% \item[]
		% \end{compactdesc}
		
% ----------------------
\subsection{Key Phrases} % *KEYPHRASES *PHRASES
% ----------------------
	\label{BOSOM-PHRASES}

	\begin{compactdesc}
	
		\item[bosom of eternity] \textbf{1} | \hyperref[DC88-13]{DC 88:13}
			\begin{compactdesc}
				\item[Bible anchor verses] ---
					% \begin{compactdesc}
						% \item[Romans 5:5] \gr{}
					% \end{compactdesc}
				\item[Commentary] There's also a use in a note to JS-H in the scriptures. It could appear elsewhere.
			\end{compactdesc}
			
	\end{compactdesc}
	
% % ----------------------
% \subsection{Bible Dictionaries} % *DICTIONARIES
% % ----------------------
	% \label{BOSOM-BD}

% % ----------------------
% \subsection{Restoration Usage} % *RESTORATION
% % ----------------------
	% \label{BOSOM-RESTORATION}

	% % -----
	% \subsubsection{Joseph Smith} % *JS
	% % -----
		% \label{BOSOM-JS}
	
		% \begin{compactdesc}
			% \item[{\hyperref[KEY]{Date/Source}}]
		% \end{compactdesc}
		
	% % -----
	% \subsubsection{Temple Ordinances}
	% % -----
		% \label{BOSOM-TEMPLE}
	
		% \begin{compactdesc}
			% \item[{\hyperref[KEY]{Date/Source}}]
		% \end{compactdesc}
	
	% % -----
	% \subsubsection{Brigham Young} % *BY
	% % -----
		% \label{BOSOM-BY}
	
		% \begin{compactdesc}
			% \item[{\hyperref[KEY]{Date/Source}}]
		% \end{compactdesc}
	
% % ----------------------
% \subsection{Commentary and Studies} % *COMMENTARY *STUDIES
% % ----------------------
	% \label{BOSOM-COMMENTARY}

	% % -----
	% \subsubsection{Key Points}
	% % -----
		% \label{BOSOM-KEYS}
	
		% \begin{compactitem}
			% \item
		% \end{compactitem}
		
	% % -----
	% \subsubsection{Sample Study}
	% % -----
		% \label{BOSOM-study}